\chapter{Peripheral Channel Structure and Transfer-Operator Gaps: Supporting Results}
\label{ch:spectral_supporting_results}

This appendix supports Chapter~\ref{ch:spectral}. It proves the algebraic
mixed-transfer identities, an auxiliary trace-positivity result, the
peripheral intertwiners behind the gauge-rigidity theorems, the Banach-algebra
convergence facts behind the overlap-decay theorems, the rank-one Perron
projection behind the primitive overlap limit, and the cyclic-decomposition
machinery that removes peripheral periodicity.

\section{Mixed-transfer powers and overlap traces}
\label{sec:app_spectral_mixed_trace}

This section proves the self-transfer identity, the word expansion of the
mixed transfer operator, and the trace identities used in
Theorem~\ref{thm:overlap_trace} and Theorem~\ref{thm:overlap_trace_rect}.
Recall the mixed transfer operator $F_{AB}$ of
Definition~\ref{def:mixed_transfer} and its rectangular form
$F_{AB}^{\mathrm{rect}}$ of Definition~\ref{def:mixed_transfer_rect}.

\begin{theorem}[Self-transfer]\label{thm:mixed_self}
    \lean{Kraus.mixedMapLM_self}
    \leanok
    \uses{def:mixed_transfer, def:transfer_map}
    Setting $B=A$ recovers the ordinary transfer map: $F_{AA}=\E_A$.
\end{theorem}

\begin{proof}\leanok
    Substituting $B=A$ into (\ref{eq:spectral_mixed_transfer}) gives
    $F_{AA}(X)=\sum_iA^iX(A^i)^\dagger=\E_A(X)$ by
    (\ref{eq:mps_transfer_map}).
\end{proof}

\begin{lemma}[Iterated mixed transfer]\label{thm:mixed_pow}
    \lean{Kraus.mixedMapLM_pow_apply}
    \leanok
    \uses{def:mixed_transfer, def:eval_word}
    For every $X\in\MN{D}$ and $N\geq0$,
    \begin{align}
        F_{AB}^N(X)
         & =\sum_{\sigma\in\{0,\ldots,d{-}1\}^N}
        A^\sigma X(B^\sigma)^\dagger,
        \label{eq:spectral_mixed_power}
    \end{align}
    where $A^\sigma=A^{\sigma_1}\cdots A^{\sigma_N}$ denotes the word
    evaluation of Definition~\ref{def:eval_word}.
\end{lemma}

\begin{proof}
    \leanok
    Induct on $N$. The base case is $F_{AB}^0(X)=X$. For the inductive
    step, apply $F_{AB}$ to each summand in
    (\ref{eq:spectral_mixed_power}) and reindex using
    $\{0,\ldots,d{-}1\}^{N+1}\cong
        \{0,\ldots,d{-}1\}\times\{0,\ldots,d{-}1\}^N$.
\end{proof}


Theorem~\ref{thm:overlap_trace_rect} combines Lemma~\ref{thm:mixed_pow} with
Lemma~\cite[``Trace expansion over matrix units'']{QICLean2026Blueprint} over the
rectangular matrix-unit basis: expanding $\Tr(T)$ for
$T=(F_{AB}^{\mathrm{rect}})^N$ by the trace-expansion formula
of~\cite[``Trace expansion over matrix units'']{QICLean2026Blueprint}, applying
(\ref{eq:spectral_mixed_power}) to each $T(E_{pq})$, and reassembling the sums
over matrix-unit indices gives
\begin{align}
    \Tr(T)
     & =\sum_\sigma
    \tr(A^\sigma)\overline{\tr(B^\sigma)}
    =\sum_\sigma
    V^{(N)}(A)_\sigma\overline{V^{(N)}(B)_\sigma}
    =O_{AB}(N).
    \notag
\end{align}
Theorem~\ref{thm:overlap_trace} is the square case $D_1=D_2=D$ of this
statement, so it needs no separate expansion.

\section{Auxiliary trace positivity}


\section{Rectangular spectral-radius formulation}

\begin{lemma}[Rectangular spectral radius bound]\label{thm:spectral_radius_le_one_rect}
    \lean{Kraus.spectralRadius_mixedTransfer₂_le_one}
    \leanok
    \uses{def:mixed_transfer_rect}
    For normalized tensors of bond dimensions $D_1$ and $D_2$,
    $\rho_{\spec}(F_{AB}^{\mathrm{rect}})\leq1$.
\end{lemma}

\begin{proof}\leanok
    This is the MPS formulation of
    the corresponding theorem in the companion quantum-channel volume~\cite[``Spectral-radius bound for a trace-preserving mixed Kraus map'']{QICLean2026Blueprint}: the rectangular
    mixed transfer is the mixed Kraus map, and the two tensor normalization
    identities are its trace-preserving hypotheses.
\end{proof}

\section{Peripheral intertwiners and gauge rigidity}
\label{sec:app_spectral_rigidity}

This section proves the intertwining relations behind
Theorem~\ref{thm:modulus_one_gauge_irr_tp} and
Theorem~\ref{thm:transfer_operator_gap_rect_irr_tp}: transporting a
modulus-one mixed-transfer eigenvector through separate unital gauges
produces Kraus-level intertwiners, and these intertwiners force equal bond
dimension once they are invertible.

\begin{lemma}[Gauged peripheral eigenvector yields Kraus intertwining]%
    \label{thm:gauged_intertwining_core}
    \lean{Kraus.gauged_intertwining_core}
    \leanok
    \uses{def:mixed_transfer_rect, thm:qpf}
    Let $A$ be a normalized tensor of bond dimension $D_1$ and $B$ a
    normalized tensor of bond dimension $D_2$. Suppose $\rho_A,\rho_B$ are
    fixed points of the respective transfer maps and $S_A,S_B$ are
    invertible matrices with $S_AS_A^\dagger=\rho_A$ and
    $S_BS_B^\dagger=\rho_B$ (for instance the positive definite
    Perron--Frobenius fixed points of Theorem~\ref{thm:qpf} and their
    square roots, as in the right-canonical gauge of
    Theorem~\cite[``Right-canonical (unital) gauge'']{QICLean2026Blueprint}). Gauge to unital families
    $A'^i=S_A^{-1}A^iS_A$ and $B'^i=S_B^{-1}B^iS_B$. If
    $X\in M_{D_1\times D_2}(\C)\setminus\{0\}$ satisfies
    $F_{AB}^{\mathrm{rect}}(X)=\mu X$ with $|\mu|=1$, then
    $X'=S_A^{-1}X(S_B^\dagger)^{-1}$ is nonzero, the families $A'$ and $B'$
    are unital, and, for every $i$,
    \begin{align}
        X'(B'^i)^\dagger
         & =\mu(A'^i)^\dagger X',
         &
        A'^iX'
         & =\mu X'B'^i.
        \label{eq:spectral_gauged_intertwining}
    \end{align}
\end{lemma}

\begin{proof}\leanok
    Embed the gauged families and transported eigenvector into the block
    matrices
    \begin{align}
        C^i
         & :=
        \begin{pmatrix}
            A'^i & 0    \\
            0    & B'^i
        \end{pmatrix},
         &
        Y
         & :=
        \begin{pmatrix}
            0 & X' \\
            0 & 0
        \end{pmatrix},
        \label{eq:spectral_block_embedding}
    \end{align}
    on $M_{D_1+D_2}(\C)$; the matrices $\{C^i\}$ form a unital Kraus family,
    and $E_C(Y)=\mu Y$. The positive definite fixed points of the adjoint
    maps give a faithful weighted trace for the block map. Applying the
    Kadison--Schwarz inequality to $Y$ and pairing its positive semidefinite
    gap with this fixed point shows that the gap vanishes. The Kraus
    relation of the corresponding theorem in the companion quantum-channel volume~\cite[``KS equality implies Kraus commutation'']{QICLean2026Blueprint} then gives
    $Y(C^i)^\dagger=\mu(C^i)^\dagger Y$; applying the same argument to the
    complementary product gives $C^iY=\mu YC^i$. Expanding these block
    relations gives (\ref{eq:spectral_gauged_intertwining}).
\end{proof}

\begin{lemma}[Intertwining yields gauge-phase equivalence]%
    \label{thm:gauged_intertwining_gauge_phase}
    \lean{MPSTensor.gaugePhaseEquiv_of_gauged_intertwining}
    \leanok
    \uses{def:gauge_phase_equiv}
    Let $A$ and $B$ be tensors of the same bond dimension. Suppose invertible
    gauges take them to $A'$ and $B'$, and suppose there are an invertible
    matrix $X'$ and a scalar $\mu$ with $|\mu|=1$ such that
    $A'^iX'=\mu X'B'^i$ for every physical index $i$. Then $A$ and $B$ are
    gauge-phase equivalent.
\end{lemma}

\begin{proof}\leanok
    If the bond dimension is zero, take the unit scalar and identity gauge;
    the matrix relation is vacuous because the bond index is empty. Otherwise,
    install the resulting nonzero-dimension hypothesis and apply the companion
    quantum-channel gauged-intertwining theorem, which combines the two gauges
    with the intertwiner into a Kraus gauge-phase witness. Finally, the local
    conversion theorem translates that witness, including its relation
    orientation, to gauge-phase equivalence of the original tensors.
\end{proof}

Combining these two lemmas proves Theorem~\ref{thm:modulus_one_gauge_irr_tp}:
Lemma~\ref{thm:gauged_intertwining_core} supplies the intertwining relations
(\ref{eq:spectral_gauged_intertwining}) for the separately gauged, unital
families. The Kadison--Schwarz equality used inside its own proof makes
$X'X'^\dagger$ and $X'^\dagger X'$ nonzero positive semidefinite fixed points
of the gauged transfer maps; under irreducibility,
Theorem~\ref{thm:transfer_psd_fp_pd_irred} upgrades these to positive definite
matrices, so the transported intertwiner $X'$ is invertible.
Lemma~\ref{thm:gauged_intertwining_gauge_phase} then converts the invertible
intertwiner, together with the two gauge matrices, into a single invertible
matrix realizing the gauge-phase relation
$B^i=\overline\mu XA^iX^{-1}$ via the Skolem--Noether theorem
(Theorem~\cite[``Skolem--Noether'']{QICLean2026Blueprint}).

Lemma~\ref{thm:rectangular_intertwining_dim_eq} gives the companion
kernel-invariance argument, independent of irreducibility: if a modulus-one
rectangular intertwiner (\ref{eq:spectral_rect_intertwining}) exists between
injective tensors of bond dimensions $D_1$ and $D_2$, then the first
relation shows that $\ker(X)$ is invariant under every $(B^i)^\dagger$;
since the $\{B^i\}$ span the full matrix algebra, so do the
$\{(B^i)^\dagger\}$, so $\ker(X)$ is trivial and $D_2\leq D_1$. Taking
adjoints of the second relation and applying the same argument to
$X^\dagger$ shows that $\ker(X^\dagger)$ is invariant under every
$(A^i)^\dagger$ and hence trivial, giving $D_1\leq D_2$. Combined with
Lemma~\ref{thm:gauged_intertwining_core}, this proves
$\rho_{\spec}(F_{AB}^{\mathrm{rect}})<1$ whenever $A$ and $B$ are injective
tensors of unequal bond dimension, which is exactly
Theorem~\ref{thm:transfer_operator_gap_rect}.

\section{Spectral-radius decay and overlap limits}
\label{sec:app_spectral_decay}

This section collects the Banach-algebra convergence facts underlying
Theorem~\ref{thm:transfer_pow_zero}, Theorem~\ref{thm:overlap_decay}, and
the block-separation consequences of the transfer-operator gap.


Theorem~\ref{thm:rectangular_transfer_pow_zero} applies
Lemma~\cite[``Powers vanish below unit spectral radius'']{QICLean2026Blueprint} to the rectangular mixed transfer
operator and then evaluates its powers at a fixed matrix.
Theorem~\ref{thm:transfer_pow_zero} combines the gap
$\rho_{\spec}(F_{AB})<1$ of Theorem~\ref{thm:transfer_operator_gap} with the
equal-bond-dimension specialization of this rectangular decay theorem.
Since $\Tr$ is continuous on the finite-dimensional matrix-endomorphism
algebra, this gives $\Tr(F_{AB}^N)\to0$ term by term over the finitely many
matrix-unit entries in the trace-expansion formula
of~\cite[``Trace expansion over matrix units'']{QICLean2026Blueprint}; combined with
the overlap--trace identity (\ref{eq:spectral_overlap_trace}), this proves
Theorem~\ref{thm:overlap_decay}.

\begin{theorem}[Cross-correlation decay]\label{thm:cross_decay}
    \lean{MPSTensor.cross_correlation_tendsto_zero}
    \leanok
    \uses{def:mixed_transfer, def:gauge_phase_equiv, def:injective}
    If $A$ and $B$ are injective normalized tensors that are not gauge-phase
    equivalent, then $\tr(F_{AB}^N(X))\to0$ as $N\to\infty$ for every
    $X\in\MN{D}$.
\end{theorem}

\begin{proof}\leanok
    \uses{thm:transfer_pow_zero}
    Theorem~\ref{thm:transfer_pow_zero} gives $F_{AB}^N(X)\to0$ in operator
    norm, and continuity of the trace gives the conclusion.
\end{proof}

\begin{theorem}[Self-correlation persists]\label{thm:self_persist}
    \lean{Kraus.self_correlation_persists}
    \leanok
    \uses{def:transfer_map}
    If $\rho$ is a fixed point of $\E_A$, then
    $\tr(\E_A^N(\rho))=\tr(\rho)$ for every $N\geq0$. In particular, this
    applies to the distinguished Perron--Frobenius fixed point from
    Theorem~\ref{thm:qpf}.
\end{theorem}

\begin{proof}\leanok
    Since $\E_A(\rho)=\rho$, one has $\E_A^N(\rho)=\rho$ for every $N$.
    Taking the trace proves the claim.
\end{proof}

\section{Rank-one Perron projection}
\label{sec:app_spectral_perron_projection}

This section proves the rank-one Perron limit
(Theorem~\ref{thm:trace_pow_one}) underlying the primitive overlap
convergence (Theorem~\ref{thm:primitive_overlap}). Let $E:\MN{D}\to\MN{D}$
be trace-preserving with a fixed point $\rho\neq0$, $\tr(\rho)\neq0$. Write
$P:=P_\rho$ for the fixed-point projection of
the corresponding definition in the companion quantum-channel volume~\cite[``Fixed-point projection'']{QICLean2026Blueprint},
\begin{align}
    P(X)
     & =\frac{\tr(X)}{\tr(\rho)}\rho,
    \notag
\end{align}
and $N:=E-P$ for the complementary part. The power decomposition
(the corresponding theorem in the companion quantum-channel volume~\cite[``Power decomposition'']{QICLean2026Blueprint}) gives
$E^n=P+N^n$ for $n\geq1$.


Under the hypothesis $\rho_{\spec}(N)<1$ of Theorem~\ref{thm:trace_pow_one},
Lemma~\cite[``Powers vanish below unit spectral radius'']{QICLean2026Blueprint} gives $N^n\to0$, so
$\Tr(E^n)=\Tr(P)+\Tr(N^n)\to\Tr(P)=1$ by
Lemma~\cite[``Trace of the fixed-point projection'']{QICLean2026Blueprint}. This is exactly the mechanism recorded
in Theorem~\ref{thm:trace_pow_one}.

When $E=\E_A$ is the transfer map of a normalized MPS tensor and $\rho$ is
its Perron--Frobenius fixed point (Theorem~\ref{thm:qpf},
Chapter~\ref{ch:qpf}), the overlap--trace identity
(\ref{eq:spectral_overlap_trace}) with $A=B$ gives
$O_{AA}(N)=\Tr(\E_A^N)$, so Theorem~\ref{thm:trace_pow_one} yields the
primitive self-overlap limit of Theorem~\ref{thm:primitive_overlap}. The
hypothesis $\rho_{\spec}(N)<1$ is exactly the complementary transfer-map gap
that the complementary spectral-radius theorem for primitive maps
(the corresponding theorem in the companion quantum-channel volume~\cite[``Complementary transfer-map gap for primitive maps'']{QICLean2026Blueprint})
supplies: for a primitive $E$ all eigenvalues bounded by $1$ in modulus and
no trace-zero fixed point, every eigenvalue of $N=E-P$ has modulus strictly
less than $1$, and since $\MN{D}$ is finite-dimensional this bounds the
spectral radius $\rho_{\spec}(N)$ itself.


\begin{theorem}[Complementary spectral-radius gap for irreducible primitive tensors]%
    \label{thm:spectralRadius_compl_lt_one_peripheralPrimitive_irred}
    \lean{Kraus.spectralRadius_compl_lt_one_of_peripheralPrimitive_of_irreducible}
    \leanok
    \uses{def:irreducible_tensor}
    Let $A$ be an irreducible normalized MPS tensor whose transfer map
    $\E_A$ is primitive. Then there is a nonzero positive semidefinite
    fixed point $\rho$ of $\E_A$ with $\tr(\rho)\neq0$ such that, for
    $P:=P_\rho$ and $N:=\E_A-P$, $\rho_{\spec}(N)<1$.
\end{theorem}

\begin{proof}\leanok
    \uses{thm:psd_fp_exists}
    The channel fixed-point existence theorem
    (Theorem~\ref{thm:psd_fp_exists}) supplies a nonzero positive
    semidefinite fixed point $\rho$. Apply
    Theorem~\cite[``Complementary gap at a prescribed fixed point'']{QICLean2026Blueprint} to
    $\E_A$ and $\rho$.
\end{proof}

This concretely instantiates the hypothesis of
Theorem~\ref{thm:trace_pow_one} and, via Theorem~\ref{thm:primitive_overlap},
gives the primitive self-overlap limit for every irreducible peripherally
primitive normalized MPS tensor.

\section{Periodicity removal}
\label{sec:app_spectral_periodicity}

This section proves the cyclic-decomposition machinery of
Theorem~\cite[``Cyclic decomposition for irreducible finite Kraus maps'']{QICLean2026Blueprint} and the
divisibility statements of Theorem~\cite[``Transfer-map peripheral eigenvalues form a cyclic group'']{QICLean2026Blueprint} and
Theorem~\cite[``Transfer-map period divides bond dimension'']{QICLean2026Blueprint}: normalizing a peripheral
eigenvector to a unitary, forming its discrete Fourier (cyclic) spectral
projections, and tracking how the channel shifts the resulting corners.


Combining these lemmas proves Theorem~\cite[``Cyclic decomposition for irreducible finite Kraus maps'']{QICLean2026Blueprint}:
Lemma~\cite[``Normalized peripheral unitary'']{QICLean2026Blueprint} supplies a unitary generator of
exact order $m$ for a primitive $m$-th peripheral root, and
Lemma~\cite[``Cyclic projections from a peripheral unitary'']{QICLean2026Blueprint} builds the $m$
Fourier projections out of its powers, using
Lemma~\cite[``Powers on a peripheral-unitary orbit'']{QICLean2026Blueprint} to identify how the channel shifts
them.

\subsection{Cyclic corners and restricted primitivity}

Corner preservation, the associated corner subspace, irreducibility of the
restricted map, and corner rank are defined in the companion quantum-channel
volume~\cite[``Corner preservation'' and the following corner definitions]{QICLean2026Blueprint}.
We use those definitions here in the MPS periodicity-removal argument.


\subsection{Group structure and divisibility}


Theorem~\cite[``Peripheral eigenspaces are one-dimensional'']{QICLean2026Blueprint} follows the same pattern:
given a peripheral unitary eigenvector $U$ for $\gamma$
(Lemma~\cite[``Peripheral unitary eigenvector'']{QICLean2026Blueprint}), the multiplicative-domain
identity gives $E(XU^\dagger)=XU^\dagger$ for any $\gamma$-eigenvector $X$,
so the scalar fixed-point lemma for irreducible unital maps
(Lemma~\cite[``Scalar fixed points for irreducible unital maps'']{QICLean2026Blueprint}) gives
$XU^\dagger=c\Id$, hence $X=cU$ and the eigenspace is one-dimensional.

The divisibility statements of Theorem~\cite[``Transfer-map peripheral eigenvalues form a cyclic group'']{QICLean2026Blueprint} and
Theorem~\cite[``Transfer-map period divides bond dimension'']{QICLean2026Blueprint} follow from
Lemma~\cite[``Cyclic projections from a peripheral unitary'']{QICLean2026Blueprint} and
Theorem~\cite[``Rank equals trace for an orthogonal projection'']{QICLean2026Blueprint}: the $m$ cyclic Fourier projections
$P_0,\ldots,P_{m-1}$ sum to $\Id$ and are cyclically permuted by $E$, so
trace preservation gives $\tr(P_{k+1})=\tr(P_k)$ for every $k$, hence
$D=\tr(\Id)=\sum_{k=0}^{m-1}\tr(P_k)=m\,\tr(P_0)$. Since the trace of an
orthogonal projection is a natural number (Theorem~\cite[``Rank equals trace for an orthogonal projection'']{QICLean2026Blueprint}),
$m\mid D$.

This periodicity-removal chain --- unitary normalization, Fourier
projections, cyclic Kraus shift (Lemma~\ref{lem:kraus_cyclic_shift}), the
blocked cyclic corners of this section, and restricted primitivity --- is
exactly the machinery behind the period-removing blocking theorems of
Chapter~\ref{ch:canonical}: \cite[Theorem~6.6]{Wolf2012Quantum} supplies the
cyclic peripheral spectrum of the adjoint transfer map, and
Theorem~\ref{thm:cyclic_sector_decomp_after_blocking} (Section~\ref{sec:cyclic_sectors})
and Theorem~\ref{thm:cyclic_sector_decomp_after_blocking_isometry} apply it
to realize the cyclic-sector decomposition of a blocked irreducible tensor
with rectangular support isometries.
